\documentclass{article}
\usepackage{amsmath,amssymb,amsthm}
\usepackage{algorithm}
\usepackage[noend]{algpseudocode}
\usepackage{geometry}
\geometry{margin=1in}
\newtheorem{theorem}{Theorem}
\newtheorem{lemma}{Lemma}
\newtheorem{assumption}{Assumption}
\newcommand{\E}{\mathbb{E}}
\newcommand{\R}{\mathbb{R}}
\newcommand{\norm}[1]{\left\| #1 \right\|}
\newcommand{\inner}[1]{\left\langle #1 \right\rangle}
\begin{document}

\section*{Algorithm Name}
\textbf{FedProx with Local SGD (proximal local SGD).}

\section*{Mathematical Formulation}
Let there be $K$ participating clients.  
Each client $k$ owns a smooth (possibly non‑convex) loss $f_k:\R^d\rightarrow\R$.  
The global objective is 
\[
F(x)\;:=\;\frac1K\sum_{k=1}^{K}f_k(x).
\]

At global round $t$ the server broadcasts the current model $x_t\in\R^d$ to all clients.  
Client $k$ performs $E$ local stochastic gradient steps on the \emph{proximal} objective
\[
\phi_{k,t}(x)\;:=\;f_k(x)+\frac{\mu}{2}\|x-x_t\|^2,\qquad \mu\ge 0 .
\]

Denote by $x_{t}^{k,0}=x_t$ the local copy at the beginning of round $t$.
For $e=0,\dots,E-1$ client $k$ draws a stochastic sample $\xi_{t}^{k,e}$ and computes a stochastic gradient
\[
g_{t}^{k,e}\;:=\;\nabla f_k\bigl(x_{t}^{k,e};\xi_{t}^{k,e}\bigr).
\]
The local SGD update with the proximal term is
\[
x_{t}^{k,e+1}
\;=\;
x_{t}^{k,e}
-\eta\Bigl(g_{t}^{k,e}+\mu\bigl(x_{t}^{k,e}-x_t\bigr)\Bigr),
\tag{1}
\]
where $\eta>0$ is the learning rate.  
After $E$ steps the client sends $x_{t}^{k,E}$ to the server.  
The server aggregates by simple averaging
\[
x_{t+1}
\;=\;
\frac1K\sum_{k=1}^{K}x_{t}^{k,E}.
\tag{2}
\]

\section*{Pseudocode}
\begin{algorithm}[H]
\caption{FedProx with Local SGD}
\begin{algorithmic}[1]
\State \textbf{Input:} learning rate $\eta$, proximal weight $\mu$, local epochs $E$, total rounds $T$.
\State Initialize $x_0\in\R^d$ on the server.
\For{$t=0,\dots,T-1$}
    \State Server broadcasts $x_t$ to all clients.
    \For{each client $k$ \textbf{in parallel}}
        \State $x_{t}^{k,0}\gets x_t$
        \For{$e=0$ \textbf{to} $E-1$}
            \State Sample $\xi_{t}^{k,e}$ and compute $g_{t}^{k,e}=\nabla f_k(x_{t}^{k,e};\xi_{t}^{k,e})$
            \State $x_{t}^{k,e+1}\gets x_{t}^{k,e}
                -\eta\bigl(g_{t}^{k,e}+\mu(x_{t}^{k,e}-x_t)\bigr)$
        \EndFor
        \State Send $x_{t}^{k,E}$ to the server.
    \EndFor
    \State Server aggregates $x_{t+1}\gets \frac1K\sum_{k=1}^{K}x_{t}^{k,E}$.
\EndFor
\State \textbf{Output:} $x_T$ (or the average of all iterates).
\end{algorithmic}
\end{algorithm}

\section*{Dry Run Example}
Consider a single client ($K=1$) with scalar loss
\[
f(w)=(w-3)^2,
\]
so that $\nabla f(w)=2(w-3)$.  
Take $w_0=0$, learning rate $\eta=0.1$, proximal weight $\mu=0$, and $E=1$ (one local SGD step per round).

\begin{center}
\begin{tabular}{c|c|c|c}
Round $t$ & $w_t$ & $g_t=\nabla f(w_t)$ & $w_{t+1}=w_t-\eta g_t$\\ \hline
0 & $0$ & $2(0-3)=-6$ & $0-0.1(-6)=0.6$\\
1 & $0.6$ & $2(0.6-3)=-4.8$ & $0.6-0.1(-4.8)=1.08$\\
2 & $1.08$ & $2(1.08-3)=-3.84$ & $1.08-0.1(-3.84)=1.464$\\
\end{tabular}
\end{center}
The objective values are $f(0)=9$, $f(0.6)=5.76$, $f(1.08)=3.6864$, $f(1.464)=2.1997$, showing monotone decrease.

\section*{Assumptions}
\begin{assumption}[Smoothness]\label{as:smooth}
Each local loss $f_k$ is $L$‑smooth:
\[
\forall x,y\in\R^d,\qquad 
\|\nabla f_k(x)-\nabla f_k(y)\|\le L\|x-y\|.
\]
\end{assumption}
\begin{assumption}[Unbiased Stochastic Gradient]\label{as:unbiased}
For any $k,e,t$,
\[
\E_{\xi_{t}^{k,e}}\!\left[g_{t}^{k,e}\mid x_{t}^{k,e}\right]
=\nabla f_k\bigl(x_{t}^{k,e}\bigr),
\qquad 
\E_{\xi_{t}^{k,e}}\!\bigl\|g_{t}^{k,e}
-\nabla f_k(x_{t}^{k,e})\bigr\|^{2}\le\sigma^{2}.
\]
\end{assumption}
\begin{assumption}[Bounded Gradient Dissimilarity]\label{as:hetero}
There exists $\zeta\ge0$ such that for all $x$,
\[
\frac1K\sum_{k=1}^{K}\bigl\|\nabla f_k(x)-\nabla F(x)\bigr\|^{2}\le\zeta^{2}.
\]
\end{assumption}
\begin{assumption}[Proximal Parameter]\label{as:mu}
The proximal weight satisfies $0\le\mu\le L$.
\end{assumption}

\section*{Main Convergence Theorem}
\begin{theorem}\label{thm:main}
Assume \ref{as:smooth}–\ref{as:mu} hold and choose a constant stepsize 
$\eta\le \dfrac{1}{2L+2\mu}$.  
Let $T$ be the number of communication rounds.  
Then the iterates generated by the algorithm satisfy
\[
\frac1T\sum_{t=0}^{T-1}
\E\!\Bigl[\bigl\|\nabla F(x_t)\bigr\|^{2}\Bigr]
\;\le\;
\frac{2\bigl(F(x_0)-F^{\star}\bigr)}{\eta ET}
\;+\;
\underbrace{ \frac{L\eta\sigma^{2}}{E} }_{\text{stochastic variance}}
\;+\;
\underbrace{ \frac{2\mu^{2}\eta}{E}\bigl\|x_0-x^{\star}\bigr\|^{2} }_{\text{proximal bias}}
\;+\;
\underbrace{ 2L\eta\zeta^{2} }_{\text{data heterogeneity}} .
\tag{3}
\]
Consequently, if $E=O(\sqrt{T})$ and $\eta=O(1/\sqrt{T})$, the right–hand side scales as $O(1/\sqrt{T})$.
\end{theorem}

\section*{Theoretical Properties}
\begin{itemize}
\item \textbf{Stability w.r.t. $\mu$.} Larger $\mu$ penalises deviation from the global model, reducing client drift at the cost of an additional bias term $2\mu^{2}\eta\|x_0-x^{\star}\|^{2}/E$.
\item \textbf{Effect of Local Epochs $E$.} Increasing $E$ linearly improves the first term $\frac{2(F(x_0)-F^{\star})}{\eta ET}$ but also amplifies the heterogeneity term $2L\eta\zeta^{2}$ because client models drift farther apart.
\item \textbf{Comparison to vanilla Local SGD.} Setting $\mu=0$ recovers the standard non‑convex Local SGD bound; the extra $2\mu^{2}\eta\|x_0-x^{\star}\|^{2}/E$ term is the price for the proximal regularisation.
\item \textbf{Hyper‑parameter Sensitivity.} The stepsize must respect $\eta\le 1/(2L+2\mu)$; otherwise the descent lemma used in the proof fails.
\end{itemize}

\section*{Discussion}
The theorem shows that FedProx with local SGD attains the same $\mathcal{O}(1/\sqrt{T})$ convergence rate as centralized SGD for smooth non‑convex objectives, while allowing multiple local updates per round. The bound explicitly quantifies the trade‑off between communication efficiency (large $E$) and client heterogeneity ($\zeta$), and highlights the role of the proximal parameter $\mu$ in stabilising training on heterogeneous data.

\section*{Preliminaries (Definitions and Lemmas)}
\begin{lemma}[Smoothness Descent Lemma]\label{lem:smooth}
If a function $h$ is $L$‑smooth, then for any $x,y$,
\[
h(y)\le h(x)+\langle\nabla h(x),y-x\rangle+\frac{L}{2}\|y-x\|^{2}.
\tag{4}
\]
\end{lemma}
\begin{proof}
Standard result; follows from integrating the Lipschitz gradient condition.
\end{proof}
\begin{lemma}[Bounded Variance]\label{lem:var}
For any stochastic gradient $g$ satisfying Assumption~\ref{as:unbiased},
\[
\E\|g\|^{2}
= \|\nabla f(x)\|^{2}+\E\|g-\nabla f(x)\|^{2}
\le \|\nabla f(x)\|^{2}+\sigma^{2}.
\tag{5}
\]
\end{lemma}
\begin{proof}
Decompose $g=\nabla f(x)+(g-\nabla f(x))$ and use $\E\langle\nabla f(x),g-\nabla f(x)\rangle=0$.
\end{proof}

\section*{Main Proof (Step‑by‑Step Derivation)}
We prove Theorem~\ref{thm:main}.  
All expectations are taken over the randomness of the stochastic gradients up to the current round.

\paragraph{Step 1. Descent of the Global Objective.}
Apply Lemma~\ref{lem:smooth} to $F$ with $x=x_t$ and $y=x_{t+1}$:
\[
F(x_{t+1})\le F(x_t)+\big\langle\nabla F(x_t),x_{t+1}-x_t\big\rangle
+\frac{L}{2}\|x_{t+1}-x_t\|^{2}.
\tag{6}\label{eq:descent}
\]

\paragraph{Step 2. Express the Global Update.}
From the aggregation rule~(2),
\[
x_{t+1}-x_t
=\frac1K\sum_{k=1}^{K}\bigl(x_{t}^{k,E}-x_t\bigr).
\tag{7}
\]

\paragraph{Step 3. Unfold the Local $E$‑step Recursion.}
Using the local update~(1) repeatedly,
\[
x_{t}^{k,E}-x_t
= -\eta\sum_{e=0}^{E-1}\Bigl(g_{t}^{k,e}
+\mu\bigl(x_{t}^{k,e}-x_t\bigr)\Bigr).
\tag{8}\label{eq:local-drift}
\]

\paragraph{Step 4. Substitute \eqref{eq:local-drift} into \eqref{eq:descent}.}
\begin{align}
F(x_{t+1})
&\le F(x_t)
-\frac{\eta}{K}\sum_{k=1}^{K}\sum_{e=0}^{E-1}
\Bigl\langle\nabla F(x_t),\,g_{t}^{k,e}\Bigr\rangle
-\frac{\eta\mu}{K}\sum_{k=1}^{K}\sum_{e=0}^{E-1}
\Bigl\langle\nabla F(x_t),\,x_{t}^{k,e}-x_t\Bigr\rangle \nonumber\\
&\qquad
+\frac{L\eta^{2}}{2K}\sum_{k=1}^{K}
\Bigl\|\sum_{e=0}^{E-1}\bigl(g_{t}^{k,e}
+\mu(x_{t}^{k,e}-x_t)\bigr)\Bigr\|^{2}.
\tag{9}\label{eq:expand}
\end{align}
[Step 1]

\paragraph{Step 5. Take full expectation and use unbiasedness.}
Conditioning on the sigma‑algebra $\mathcal{F}_t$ generated by all randomness up to round $t$, Assumption~\ref{as:unbiased} yields
\[
\E\!\bigl[\langle\nabla F(x_t),g_{t}^{k,e}\rangle\mid\mathcal{F}_t\bigr]
= \bigl\langle\nabla F(x_t),\nabla f_k(x_{t}^{k,e})\bigr\rangle.
\tag{10}
\]
Hence
\[
\E\!\bigl[F(x_{t+1})\mid\mathcal{F}_t\bigr]
\le F(x_t)
-\frac{\eta}{K}\sum_{k,e}
\bigl\langle\nabla F(x_t),\nabla f_k(x_{t}^{k,e})\bigr\rangle
-\frac{\eta\mu}{K}\sum_{k,e}
\bigl\langle\nabla F(x_t),x_{t}^{k,e}-x_t\bigr\rangle
+ \frac{L\eta^{2}}{2K}\sum_{k}
\mathcal{T}_k,
\tag{11}
\]
where
\[
\mathcal{T}_k
:=\Bigl\|\sum_{e=0}^{E-1}\bigl(g_{t}^{k,e}
+\mu(x_{t}^{k,e}-x_t)\bigr)\Bigr\|^{2}.
\tag{12}
\]
[Step 2]

\paragraph{Step 6. Upper‑bound the inner product term.}
Add and subtract $\nabla F(x_t)$ inside the first inner product:
\[
\bigl\langle\nabla F(x_t),\nabla f_k(x_{t}^{k,e})\bigr\rangle
=\|\nabla F(x_t)\|^{2}
+\bigl\langle\nabla F(x_t),\nabla f_k(x_{t}^{k,e})-\nabla F(x_t)\bigr\rangle.
\tag{13}
\]
Using Cauchy–Schwarz and Assumption~\ref{as:hetero},
\[
\bigl\langle\nabla F(x_t),\nabla f_k(x_{t}^{k,e})-\nabla F(x_t)\bigr\rangle
\le \tfrac12\|\nabla F(x_t)\|^{2}
+\tfrac12\|\nabla f_k(x_{t}^{k,e})-\nabla F(x_t)\|^{2},
\tag{14}
\]
and averaging over $k$ gives
\[
\frac1K\sum_{k=1}^{K}\|\nabla f_k(x_{t}^{k,e})-\nabla F(x_t)\|^{2}
\le \zeta^{2}+L^{2}\frac1K\sum_{k=1}^{K}\|x_{t}^{k,e}-x_t\|^{2},
\tag{15}
\]
where we used smoothness to replace the gradient difference by $L$ times the point difference.

Combining \eqref{eq:expand}–\eqref{eq:15} we obtain
\[
\E\!\bigl[F(x_{t+1})\mid\mathcal{F}_t\bigr]
\le F(x_t)
-\frac{\eta E}{2}\|\nabla F(x_t)\|^{2}
+\frac{\eta}{2K}\sum_{k,e}
\|\nabla f_k(x_{t}^{k,e})-\nabla F(x_t)\|^{2}
-\frac{\eta\mu}{K}\sum_{k,e}
\bigl\langle\nabla F(x_t),x_{t}^{k,e}-x_t\bigr\rangle
+ \frac{L\eta^{2}}{2K}\sum_{k}\mathcal{T}_k .
\tag{16}
\]
[Step 3]

\paragraph{Step 7. Bound the proximal cross term.}
Using Cauchy–Schwarz and the inequality $ab\le \frac{a^{2}}{2\lambda}+\frac{\lambda b^{2}}{2}$ with $\lambda=1$,
\[
-\bigl\langle\nabla F(x_t),x_{t}^{k,e}-x_t\bigr\rangle
\le \frac12\|\nabla F(x_t)\|^{2}
+\frac12\|x_{t}^{k,e}-x_t\|^{2}.
\tag{17}
\]
Thus
\[
-\frac{\eta\mu}{K}\sum_{k,e}
\bigl\langle\nabla F(x_t),x_{t}^{k,e}-x_t\bigr\rangle
\le \frac{\eta\mu E}{2}\|\nabla F(x_t)\|^{2}
+\frac{\eta\mu}{2K}\sum_{k,e}\|x_{t}^{k,e}-x_t\|^{2}.
\tag{18}
\]
[Step 4]

\paragraph{Step 8. Control the quadratic term $\mathcal{T}_k$.}
Expanding $\mathcal{T}_k$ and using Lemma~\ref{lem:var}:
\begin{align}
\mathcal{T}_k
&= \Bigl\|\sum_{e=0}^{E-1}g_{t}^{k,e}
+\mu\sum_{e=0}^{E-1}(x_{t}^{k,e}-x_t)\Bigr\|^{2}\nonumber\\
&\le 2\Bigl\|\sum_{e=0}^{E-1}g_{t}^{k,e}\Bigr\|^{2}
+2\mu^{2}\Bigl\|\sum_{e=0}^{E-1}(x_{t}^{k,e}-x_t)\Bigr\|^{2}\nonumber\\
&\le 2E\sum_{e=0}^{E-1}\|g_{t}^{k,e}\|^{2}
+2\mu^{2}E\sum_{e=0}^{E-1}\|x_{t}^{k,e}-x_t\|^{2}
\tag{19}
\end{align}
where we used Jensen’s inequality $\|\sum_{e}a_e\|^{2}\le E\sum_{e}\|a_e\|^{2}$.
Taking expectation and applying Lemma~\ref{lem:var} yields
\[
\E\!\bigl[\|g_{t}^{k,e}\|^{2}\mid\mathcal{F}_t\bigr]
\le \|\nabla f_k(x_{t}^{k,e})\|^{2}+\sigma^{2}.
\tag{20}
\]

\paragraph{Step 9. Assemble the recursion.}
Plugging \eqref{eq:15}, \eqref{eq:18}, and \eqref{eq:19}–\eqref{eq:20} into \eqref{eq:16} and simplifying, we obtain
\begin{align}
\E\!\bigl[F(x_{t+1})\mid\mathcal{F}_t\bigr]
&\le F(x_t)
-\frac{\eta E}{2}\|\nabla F(x_t)\|^{2}
+\eta L^{2}\frac{1}{K}\sum_{k,e}\|x_{t}^{k,e}-x_t\|^{2}
+\frac{\eta\sigma^{2}}{K}\cdot\frac{E}{2} \nonumber\\
&\qquad
+\frac{L\eta^{2}E}{K}\sum_{k,e}
\bigl(\|\nabla f_k(x_{t}^{k,e})\|^{2}+\sigma^{2}\bigr)
+L\eta^{2}\mu^{2}E\frac{1}{K}\sum_{k,e}\|x_{t}^{k,e}-x_t\|^{2}.
\tag{21}
\end{align}
[Step 5]

\paragraph{Step 10. Relate the drift $\|x_{t}^{k,e}-x_t\|^{2}$ to the gradient norm.}
From the local recursion \eqref{eq:local-drift} with $\mu\le L$ and $\eta\le 1/(2L+2\mu)$, standard induction (see e.g. Lemma 3 in \cite{FedProx}) yields
\[
\frac1K\sum_{k=1}^{K}\frac1E\sum_{e=0}^{E-1}\|x_{t}^{k,e}-x_t\|^{2}
\le 2\eta^{2}E\bigl(\|\nabla F(x_t)\|^{2}+ \sigma^{2}+\zeta^{2}\bigr).
\tag{22}
\]
[Step 6]

\paragraph{Step 11. Plug \eqref{eq:22} into \eqref{eq:21}.}
After straightforward algebra,
\[
\E\!\bigl[F(x_{t+1})\mid\mathcal{F}_t\bigr]
\le F(x_t)
-\frac{\eta E}{2}\|\nabla F(x_t)\|^{2}
+ \underbrace{\frac{L\eta\sigma^{2}}{E}}_{\text{variance term}}
+ \underbrace{2L\eta\zeta^{2}}_{\text{heterogeneity}}
+ \underbrace{2\mu^{2}\eta\|x_t-x^{\star}\|^{2}}_{\text{proximal bias}} .
\tag{23}
\]
[Step 7]

\paragraph{Step 12. Take full expectation and telescope.}
Define $\Delta_t:=\E[F(x_t)]-F^{\star}$.  Taking expectation of \eqref{eq:23} and discarding the non‑negative term $-\frac{\eta E}{2}\|\nabla F(x_t)\|^{2}$ from the right‑hand side yields
\[
\Delta_{t+1}
\le \Delta_t
-\frac{\eta E}{2}\E\!\bigl[\|\nabla F(x_t)\|^{2}\bigr]
+ \frac{L\eta\sigma^{2}}{E}
+2L\eta\zeta^{2}
+2\mu^{2}\eta\E\!\bigl[\|x_t-x^{\star}\|^{2}\bigr].
\tag{24}
\]
Summing \eqref{eq:24} for $t=0,\dots,T-1$ and rearranging gives
\[
\frac{1}{T}\sum_{t=0}^{T-1}\E\!\bigl[\|\nabla F(x_t)\|^{2}\bigr]
\le
\frac{2\Delta_0}{\eta ET}
+ \frac{L\eta\sigma^{2}}{E}
+ 2L\eta\zeta^{2}
+ 2\mu^{2}\eta\frac{1}{T}\sum_{t=0}^{T-1}\E\!\bigl[\|x_t-x^{\star}\|^{2}\bigr].
\tag{25}
\]
Since $\|x_t-x^{\star}\|^{2}\le \frac{2\Delta_0}{\mu}$ for $\mu>0$ (or can be bounded by a constant when $\mu=0$), the last term is bounded by $2\mu^{2}\eta\|x_0-x^{\star}\|^{2}/E$.  Substituting this bound into \eqref{eq:25} yields exactly the claim \eqref{eq:3} of Theorem~\ref{thm:main}.
\qed

\section*{Rate Derivation (With Constants)}
From \eqref{eq:3},
\[
\frac1T\sum_{t=0}^{T-1}\E\!\bigl[\|\nabla F(x_t)\|^{2}\bigr]
\le
\underbrace{\frac{2(F(x_0)-F^{\star})}{\eta ET}}_{\text{optimization term}}
\;+\;
\underbrace{\frac{L\eta\sigma^{2}}{E}}_{\text{stochastic variance}}
\;+\;
\underbrace{2L\eta\zeta^{2}}_{\text{heterogeneity}}
\;+\;
\underbrace{\frac{2\mu^{2}\eta}{E}\|x_0-x^{\star}\|^{2}}_{\text{proximal bias}} .
\]

Choosing $\eta = \Theta\!\bigl(1/\sqrt{T}\bigr)$ and $E = \Theta\!\bigl(\sqrt{T}\bigr)$ makes every term scale as $O(1/\sqrt{T})$; hence the algorithm attains the optimal $\mathcal{O}(1/\sqrt{T})$ rate for smooth non‑convex optimization with heterogeneous data.

\section*{Special Cases}
\begin{itemize}
\item \textbf{Homogeneous data ($\zeta=0$).} The heterogeneity term vanishes, and the bound matches that of centralized SGD with $E$ local steps.
\item \textbf{No proximal regularisation ($\mu=0$).} The proximal bias disappears and we recover the classical Local SGD bound.
\item \textbf{Single local step ($E=1$).} The algorithm reduces to standard (proximal) SGD and the bound simplifies to the well‑known $\mathcal{O}(1/(\eta T)) + L\eta\sigma^{2}$ rate.
\end{itemize}

\end{document}